\chapter{A Worked Instantiation: Reliability-Gated Multimodal Fusion (ELARA)}\label{ch:elara}

This chapter presents a fully worked, machine-checked instantiation of the K-Bound
trichotomy in the domain of multimodal anomaly detection.  The instantiation is
called \textbf{ELARA} (Evidence-Layered Anomaly Reliability Architecture) and its
meta-routing variant \textbf{ELARA-U}.  Its purpose in this monograph is not to
present a new anomaly-detection system but to show that the abstract boundary
$\{$adapt, freeze, abstain$\}$ has a concrete, fully implemented, and empirically
validated special case --- one where every theorem in the general theory maps to a
runnable code module, a validator script, and a generated artifact.  The honest
negatives in the ELARA evidence are as instructive as the positives; both are
reported without omission.

% -----------------------------------------------------------------------
\section{The Fusion and Routing Problem}\label{sec:elara-problem}
% -----------------------------------------------------------------------

\paragraph{Setting.}
Consider a heterogeneous detector zoo $\{f_1, \ldots, f_K\}$ evaluated on an
anomaly-detection task.  A labelled validation slice is available; test labels are
not.  The decision is: given the detectors' scores and the validation evidence,
should one (i) \emph{fuse} all detectors (e.g.\ stack them into a logistic ensemble),
(ii) \emph{select} the best single detector, or (iii) apply a \emph{reliability
gate} that routes to different fusion strategies depending on observable drift
signals?  This is the label-free adapt/freeze/abstain decision instantiated for
multimodal fusion.

\paragraph{Reliability-Gated Attention (RGA).}
The ELARA source system implements a reliability gate: a drift detector
(windowed KS test on detector score distributions) that routes between a
confidence-weighted mean of all detectors (the ``fused'' strategy, analogous to
$f_a$) and a fallback to the best-validation detector (analogous to $f_0$).  When
the gate fires --- indicating detectable distribution shift in one modality's
scores --- fusion is activated; when it does not, the frozen single-detector is
maintained.  This is a concrete instantiation of the K-Bound gate
(Theorem~\ref{thm:gate}): route when the evidence $Z$ certifies benefit, stay
frozen when it does not.

\paragraph{The cross-domain meta-routing benchmark (ELARA-U).}
To stress-test the gate beyond the single in-domain setting, ELARA-U evaluates the
same routing decision on a 123-task archive spanning five anomaly-detection families:
tabular (ADBench~\cite{han2022adbench}), image-OOD, text, cyber, and fraud.  Each
task provides six detector scores on labelled validation and test.  The frozen
baseline $f_0$ is the best-validation detector; the adapted candidate $f_a$ is a
rank-normalised logistic stack.  The benchmark is label-free except for the
validation slice; test labels are withheld until scoring.

% -----------------------------------------------------------------------
\section{The Verified ELARA Claim}\label{sec:elara-claim}
% -----------------------------------------------------------------------

The ELARA-U benchmark yields the following verified claim, drawn from the
repository README (ELARA-U section)~\cite{kbound2026repo}:

\begin{quote}
\textit{Stacking beats selection; reliability routing helps if and only if
modalities can fail independently.}
\end{quote}

We unpack each part of this claim in turn, reporting every number directly from
the source.

\subsection{Stacking beats selection (positive result)}\label{sec:elara-stacking}

On the 123-task / 5-family benchmark, the rank-normalised logistic stack beats
validation auto-selection by $+0.036$ AUROC (CI $[0.023,\, 0.050]$; 86/123 wins)
and the best fixed single detector by $+0.075$ (CI $[0.052,\, 0.101]$).  All
comparisons are leakage-free and Holm-corrected for multiple testing.

This result is robust:
\begin{itemize}
\item Comprehensive baselines (D28): the stack beats all nine non-oracle baselines
  (best-fixed, auto-select, MetaOD-style, rank/CW/raw/calibrated averages, AutoML
  greedy ensemble, per-family specialist) \emph{and} both test-label oracle
  ceilings (best single $+0.025$, best pair $+0.032$; all CIs exclude 0).
\item Stacker family: RF, GBM, and HistGBM match the logistic stack; the simple
  average loses.
\item Learned selection (D25): a MetaOD-style meta-selector (LOO, 28 meta-features)
  scores \emph{below} naive auto-select by $-0.011$ (CI $[-0.017,\, -0.006]$);
  the stack beats it by $+0.044$ (CI $[+0.030,\, +0.060]$).  Combining, not
  selecting, is the winning move.
\item Sealed external (D24): with the method frozen before scoring, the stack wins
  $+0.016$ on 74 untouched ADBench tasks under a different feature extractor
  (ViT/RoBERTa), CI $[+0.004,\, +0.032]$.
\end{itemize}

In K-Bound terms: stacking is the adapted candidate $f_a$, and this result establishes
the \emph{knowably helpful} regime for the fusion decision on this domain.

\subsection{Reliability routing adds nothing on single-input data (honest negative)}\label{sec:elara-routing-neg}

\textbf{Reliability/drift routing adds no value on single-input data across four
deployment regimes plus a sealed natural temporal-shift benchmark (D22).}  More
pointedly: a reliability gate \emph{significantly hurts} the stack (Holm $p < 10^{-8}$).

This is a strong negative result and we state it without softening.  On the 123-task
benchmark, activating the windowed-KS gate on top of the logistic stack causes the
stack's AUROC to drop; the Holm-corrected test rejects the null of no harm at
$p < 10^{-8}$.  The gate fires spuriously on benign mixture re-weighting even when
no real modality has failed (this is T2 in the theorem stack, quantified below), and
this spurious firing routes away from the correct fusion to the less accurate single
detector.

The practical interpretation is direct: on unimodal or single-sensor data, the
observable drift signal $Z$ lacks the specificity to distinguish helpful from harmful
regime; equivalently, the evidence is in the \emph{unknowable} region of the K-Bound
phase diagram.  Abstaining is the right action, and gating here is worse than
simply stacking.

\subsection{Routing wins where it structurally should: independent modality failure}\label{sec:elara-routing-pos}

\textbf{But routing wins where it structurally should.}  On two real multimodal
datasets with independent modality failure (D23) --- Real-IAD-D3 (15 categories,
\cite{wang2024realiad}) and MVTec-3D (7 categories,
\cite{bergmann2022mvtec3d,bergmann2019mvtec}) --- drift-gated fusion recovers
substantial performance:

\begin{itemize}
\item \textbf{Real-IAD-D3:} $0.517 \to 0.715$ AUROC under modality failure.
  Hypothesis gains: H1 $+0.241$, H2 $+0.386$, H3 $+0.248$; all CIs exclude zero.
\item \textbf{MVTec-3D:} H1 $+0.207$, H2 $+0.317$, H3 $+0.211$; all CIs exclude
  zero on both datasets.
\item On \emph{clean} data: reliability gating does no harm ($0.882 = 0.882$ AUROC,
  no change on MVTec-3D clean).
\item In a stress sweep: routing is non-significant at severity 0 ($-0.0039$,
  CI includes zero) and significant only at severity 3 ($+0.039$, CI excludes
  zero) --- the predicted knowability crossover.
\end{itemize}

This is the K-Bound phase boundary made concrete: below the drift threshold the
system is in the unknowable/helpfully-dominated region and gating costs; above the
threshold detectable modality failure places the system in the knowably-helpful
region and gating recovers.

% -----------------------------------------------------------------------
\subsection{Controlled confirmation: a significant beats-both under detectable modality failure}\label{sec:elara-controlled}
% -----------------------------------------------------------------------

The routing wins above are real but per-dataset and small-$n$.  To test the mechanism
\emph{cleanly} we ran a pre-registered controlled experiment (Protocol~D33) in which the named
condition holds by construction.  MNIST is split into two views (left/right half) as two
``sensors''; one view is degraded by Gaussian noise across $13$ severities $\times 10$ batches
$=130$ conditions.  The \textsc{adapt} candidate is late fusion of the two views; the
\textsc{freeze} default is the single clean view.  \KGA{} routes from \emph{label-free} evidence (the
noisy view's confidence and entropy, and cross-view disagreement) under a leave-one-condition-out
conformal certificate ($\alpha=0.10$).

\begin{center}
\begin{tabular}{lc}
\toprule
Policy & mean accuracy (130 conditions) \\
\midrule
always-fuse (adapt) & $0.583$ \\
always-single view (freeze) & $0.854$ \\
\KGA{}-routed & $\mathbf{0.857}$ \\
oracle & $0.857$ \\
\bottomrule
\end{tabular}
\end{center}

\KGA{} \textbf{significantly beats both} trivial policies --- $+0.274$ over always-fuse and $+0.003$
over always-single, each at $P{=}1.0$ by paired bootstrap over conditions --- with \textbf{zero
false-adapt}: it adapts on the clean-view conditions where fusion helps and freezes where the degraded
view would poison fusion.  \emph{Honest scope:} this is a \emph{controlled construction}, not a natural
benchmark.  It shows the certificate routes correctly \emph{when} the named condition holds and is
label-free detectable --- isolating the mechanism from the open empirical question of which natural
benchmarks exhibit the condition (Chapter~\ref{ch:discussion}) --- and most of its margin over the safe
single-view baseline is the safety half (avoiding fusion collapse), the gain being concentrated in the
few clean conditions.

% -----------------------------------------------------------------------
\section{The T1--T9 + GDR Theorem Stack}\label{sec:elara-theorems}
% -----------------------------------------------------------------------

The ELARA instantiation is backed by a stack of ten theorems, each with a code
module, a validator script, and a generated \LaTeX{} artifact (all 10 pass their
validators; verified by \texttt{validate\_theorem\_stack.py} as reported in
\texttt{THEOREM\_CODE\_STATUS.md} Part 2~\cite{kbound2026repo}).  Table~\ref{tab:elara-stack}
maps each theorem to its one-line statement and its role in the K-Bound general theory.

\begin{table}[ht]
\centering
\small
\caption{ELARA theorem stack (T1--T9, GDR): one-line statements and K-Bound roles.
Source: \texttt{THEOREM\_CODE\_STATUS.md} Part 2 and \texttt{kbound.tex}
Appendix~A~\cite{kbound2026repo}.}
\label{tab:elara-stack}
\begin{tabular}{p{0.05\textwidth}p{0.43\textwidth}p{0.43\textwidth}}
\toprule
ID & One-line statement & K-Bound role \\
\midrule
T1 & Reliability-blind fusion is dominated under sparse/coherent corruption
   (positive lower bound).
 & Establishes the \emph{knowably-helpful} (stress) regime: gating is
   necessary when corruption is present. Supports Theorem~\ref{thm:imp} (helpful
   side). \\
T2 & Global KS drift is inflated by benign mixture re-weighting even with no real
   drift.
 & Observability hazard: evidence $Z$ can false-fire, routing into false-adapt or the
   unknowable regime. \\
T3 & Closed-form mean-gate miss probability (CLT approximation).
 & Quantifies gate error; calibrates the certificate. Gate-error calibration for
   Theorem~\ref{thm:cert}. \\
T4 & Finite-sample risk-dominance sample complexity for the switch.
 & How much validation data certifies the helpful regime.
   \textbf{Instantiates Theorem~\ref{thm:pos}} (covariate-shift positive regime). \\
T5 & Empirical-Bernstein finite-sample switching certificate.
 & \textbf{Directly instantiates Theorem~\ref{thm:cert}}: the false-adapt-controlled
   adapt/freeze/abstain rule. \\
T6 & Windowed-KS gate as a CUSUM detector; ARL/AED trade-off.
 & Detectability of the shift trigger: when $Z$ is informative (governs the
   boundary of the unknowable regime). \\
T7 & Massart finite-class Rademacher bound for the meta-router.
 & Generalization of the benefit estimator $\widehat{\Delta}$ in
   Theorem~\ref{thm:cert}. \\
T8 & Validation-only certified selection among $\{$gate, TTA, stack$\}$.
 & Multi-candidate \KGA{} decision (extends beyond binary adapt/freeze). \\
T9 & No fusion rule beats the confidence-weighted mean on clean near-separable
   transfer.
 & \textbf{Proves the knowably-harmful/neutral (clean) regime}: freezing is
   optimal; matches Theorem~\ref{thm:imp} harmful-side. \\
GDR & Coherence-certified switch is minimax over coherent-shift vs.\
    heterogeneous-mixture adversaries.
 & \textbf{Justifies the abstain/gate action} as minimax-safe.
   Instantiates Theorem~\ref{thm:imp} (the unknowable boundary). \\
\bottomrule
\end{tabular}
\end{table}

\paragraph{How the theorems partition the trichotomy.}
T1, T3, and T4 together lower-bound the \emph{helpful} regime: they establish
conditions under which fusing reliably outperforms the frozen single detector and
certify how much validation data is needed (T4, the engine of
Theorem~\ref{thm:pos}).  T9 proves the \emph{harmful/neutral} regime: on clean
near-separable data no routing rule beats the frozen confidence-weighted mean, so
the certificate correctly refuses to adapt.  T2, T6, and T7 govern the
\emph{observability} of the evidence $Z$: T2 shows $Z$ can false-fire, T6 gives
the detection-theoretic conditions under which $Z$ is actually informative, and T7
bounds how well the benefit estimator $\widehat{\Delta}$ generalises from
validation to test.  T5 supplies the operational certificate (the engine of
Theorem~\ref{thm:cert}).  T8 extends the binary decide-or-freeze rule to a
multi-candidate selection.  GDR provides the minimax argument that certifies
abstention as the safe action when no clear signal is available (Theorem~\ref{thm:imp}).

\paragraph{Scope caveats.}
Several theorems (T2, T3, T6) are closed-form operationalizations validated
empirically, not deep proofs; this is how they are scoped in the paper and in
\texttt{THEOREM\_CODE\_STATUS.md}.  T7 has no dedicated \texttt{theory/} module
(it lives in the \texttt{audit\_meta\_router\_pac*} scripts) but is validated.
GDR has partial real-data validation.  These caveats are stated transparently;
none affects the K-Bound general theorems (Theorems~\ref{thm:imp}--\ref{thm:disagree}),
which are proved independently.

% -----------------------------------------------------------------------
\section{How ELARA Instantiates the Evidence $Z$}\label{sec:elara-z}
% -----------------------------------------------------------------------

The general K-Bound theory posits a label-free evidence vector $Z$ from which
the certificate constructs $\widehat{\Delta} \pm \varepsilon$ (Theorem~\ref{thm:cert}).
In the ELARA instantiation, $Z$ consists of:
\begin{itemize}
\item The windowed KS statistic between the current and reference score distributions
  for each modality;
\item The \emph{disagreement} between detectors (the sign-of-difference term central
  to Theorems~\ref{thm:disagree}--\ref{thm:disagree-reg});
\item Log-batch-size (governs the conformal radius $\varepsilon$);
\item The validation-slice mean performance of the fusion candidate.
\end{itemize}

The T6 theorem (KS gate as CUSUM detector) establishes the detectability conditions
under which the KS component of $Z$ is genuinely informative about the sign of the
benefit.  When KS signals are below the detection threshold the system is in the
unknowable regime (GDR), justifying abstention.

The ablation from Chapter~\ref{ch:exp1} (Table corresponding to \texttt{kbound.tex}
Table~7) confirms the relative importance of these components in the anomaly-routing
domain: removing the disagreement feature from $Z$ raises regret from $0.037$ to
$0.094$ ($2.6{\times}$), directly supporting Theorem~\ref{thm:disagree}.  The
ELARA multimodal experiments show the same pattern at the system level: the KS
gate fires appropriately on Real-IAD-D3 and MVTec-3D where modality failure is
genuine and informative, and fires spuriously on single-input data where the $Z$
signal is uninformative.

% -----------------------------------------------------------------------
\section{What Generalises and What Is Domain-Specific}\label{sec:elara-generalize}
% -----------------------------------------------------------------------

The ELARA evidence covers a specific fusion domain (multimodal anomaly
detection).  We are explicit about what in this evidence is general and what is
domain-specific.

\paragraph{What generalises.}
\begin{enumerate}
\item \textbf{The trichotomy structure.}  Helpful/harmful/unknowable regimes appear
  in ELARA exactly as predicted by the general theory: helpful under independent
  modality failure (Real-IAD-D3, MVTec-3D with corruption), harmful on clean data
  (T9: fusion is dominated), unknowable on single-input data (gating fires
  spuriously, net harm Holm $p{<}10^{-8}$).
\item \textbf{The phase boundary.}  The knowability crossover (non-significant at
  severity 0, significant at severity 3) matches the theoretical phase boundary:
  the adapted candidate's value emerges above a detectable-drift threshold.
\item \textbf{The safety guarantee.}  The T5 certificate (Theorem~\ref{thm:cert})
  delivers $0\%$ false-adapt on the 123-task benchmark harmful regime
  (\texttt{kbound\_harmful\_results.json}: \KGA{} makes 0 false-adapt decisions,
  regret $0.0019$ vs.\ always-adapt $0.0214$, an $11{\times}$ reduction).
\item \textbf{Disagreement as primary evidence.}  T7 (meta-router generalization)
  and the ablation both show disagreement is the most informative component of $Z$
  across domains, consistent with Theorem~\ref{thm:disagree}.
\end{enumerate}

\paragraph{What is domain-specific.}
\begin{enumerate}
\item \textbf{The fusion architecture.}  The logistic stack, the KS drift detector,
  and the confidence-weighted mean fallback are all specific to the anomaly-score
  stacking domain.  The theorems T1--T9 are instantiation results, not the general
  trichotomy.
\item \textbf{The 123-task helpful-dominated character.}  The cross-domain archive
  is helpful-dominated for the stacking candidate (stacking almost always beats
  selection), which differs from the CIFAR-10-C stress grid.  The certificate's
  raw-accuracy advantage is bounded by how often harmful adaptation occurs; on this
  domain it is mainly a safety certificate, not a knockout win.
\item \textbf{Real-IAD and MVTec-3D specificity.}  The modality-failure recovery
  result relies on independent modality failures that generate detectable KS drift.
  Other multimodal data with correlated modality failures would not exhibit this
  crossover.
\item \textbf{CLT/Gaussian approximations.}  T3 and parts of T6 use CLT/Gaussian
  approximations; their validity depends on regime-specific sample-size conditions
  not universally met.
\end{enumerate}

\paragraph{Summary.}
The ELARA instantiation serves two purposes in this monograph.  First, it provides
a complete, machine-verified worked example that grounds the abstract K-Bound
framework: every general theorem has a corresponding runnable code module, and the
theory's predictions (helpful/harmful/unknowable regimes, the phase crossover,
the safety certificate) are all observed in the data.  Second, the honest negative
--- routing hurts on single-input data (Holm $p{<}10^{-8}$) --- demonstrates that
the certificate is calibrated: it correctly identifies when the evidence $Z$ does
not support adaptation and recommends abstention, which in this domain means
maintaining the stack without the gate.  A framework that overclaims would not
survive this test; the fact that ELARA passes both the positive and negative halves
is evidence for the theory's calibration.
